\section{データ生成メカニズムとモデリング}
\subsection{授業内容}
\subsubsection{科目の中でのこのコマの位置づけ}
モデリングアプローチには実質的にベイズ推定が必須であり，そのためにはまず推定法としてのベイズ法の位置付け，ベイズの定理の基礎，実践的方法としてのMCMC法について理解する。特に，従来の心理統計ではモーメント法と最尤法に言及されるにとどまるものが多い。確率モデルとしての表現は最尤法から導入されており，尤度による推定を事後分布という確率分布で表現するようになったものとして，ベイズの定理を位置付けると良い。MCMC法は新しい方法であるが，その前に確率分布の特徴を記述するために乱数を利用することができる，という事実を確認すれば，事後分布の記述との相性の良さも明らかである。このように，ベイズ法を過度に新規で全く異なるものであるという印象を与えることなく，従来の方法の延長線上にあるものとして考えられるようにする。

\subsubsection{コマ主題細目}
\begin{description}
	\item[データ生成モデリング] 推測統計学が母集団分布における仮定から母数を推定することを目的とし，モーメント法，最尤法，ベイズ法といったアプローチで推定を行うものであったことを再確認する。その上で，帰無仮説検定など心理学一般で使われているモデルは得られた結果のみに基づいてモデル比較をする形である。これは客観性を重視しデータにいかなる前提も置かないことを考えてのアプローチであり，いわばデータ駆動型の分析方法である。これに対して，データがどのように生まれてきたかを考え，その仕組みに基づいて分析する方法がデータ生成モデリングである。データ駆動型分析法は実験方法にそのメカニズムを埋め込んでおり，データ生成モデル駆動型分析法はメカニズムそのものを検証する形になっている。

		\begin{flushright}
			$\to$\citet{松浦201610}のPp.18--25.
		\end{flushright}

	\item[ベイズ推定の基礎] データ生成モデル駆動型分析にあたっては，未知なるパラメータが多くなるため，実質的にベイズ推定を使うことになる。ベイズ推定の原理となる確率やベイズの公式を再確認することが必要である。

		\begin{flushright}
			$\to$\citet{Maeda2017}のPp104--123ほか枚挙に遑がない
		\end{flushright}

	\item[MCMC] ベイズ推定は事後分布を用いて分析結果を考えることになり，これは結果がパラメータの確率分布として得られることを意味する。そこで確率分布を分析する方法として，乱数を用いてアプローチすることを考える。乱数を用いるアプローチの利点は，積分計算が記述統計量で済むこと，周辺化分布についても当該変数について考えるだけで済むこと，精度を上げる時にサンプル数を増やすだけで良いことなどが挙げられる。また乱数を用いるアプローチに対応したソフトウェアとしてJAGSやStanが挙げられる。これらが確率型プログラミング言語が乱数発生機であることを踏まえ，簡単な実践を行ってみる。

		\begin{flushright}
			$\to$\citet{Maeda2017}のPp.147--194.
		\end{flushright}

		\item[乱数によるアプローチの例]簡単な確率分布を使って，乱数によるアプローチを実践する。サンプルサイズを増やすことで精度が上がること，記述統計量が確率分布の特徴を記述することを再確認する。特に平均値，中央値，パーセンタイル，分散や標準偏差など，分布の要約統計量の計算について，Rスクリプトで算出できるよう練習する。
\end{description}


\subsubsection{キーワード}
\begin{itemize}
	\item データ生成モデリング
	\item ベイズの定理
	\item マルコフ連鎖モンテカルロ法
	\item 乱数による近似
\end{itemize}
\subsection{授業情報}
\paragraph{コマの展開方法}
講義/遠隔/演習
\subsubsection{予習・復習課題}
\paragraph{予習}環境の準備が整っていないものは，急ぎ準備を行うこと。また1年次の心理学データ解析基礎において，確率分布や乱数を用いるアプローチについても言及はしているので，振り返って「確率という数字の公理」を確認しておくことが望ましい。
\paragraph{復習}Rを用いて乱数を発生させ，理論上の値の近似値になることを確認する。正規分布，ベルヌーイ分布，二項分布，ポアソン分布など複数の既存関数を確かめることで，一般化されない事後分布についての乱数発生機が手に入ったことの重要性に気づくことができるだろう。